\section{Model}
\label{sec:rewrite-model}

Time is continuous.  The economy contains a representative household, competitive final-good firms, and an integrated AI developer. Labor-augmenting technology $A$ grows exogenously. AI efficiency $B$ improves through autonomous AI research: the only current input used to raise $B$ is research compute. The model therefore describes an economy in which there are no longer human researchers developing AI and AI has achieved recursive self-improvement (RSI). As in \citet{trammellkorinek2023}, I distinguish automation of research from automation of production: human labor remains an input in final-good production.

%I derive the model in the order in which its agents solve their problems. I first present the household and final-good firms. I then state the developer's intertemporal problem before deriving its within-date service choice and its dynamic research conditions. Only after every agent has been characterized do I combine their conditions in the definition of equilibrium.

\tikzset{
 rewrite-input/.style={draw, rounded corners=2pt, minimum width=2.7cm,
  minimum height=1.0cm, align=center, fill=gray!8},
 rewrite-state/.style={draw, rounded corners=2pt, minimum width=2.7cm,
  minimum height=1.0cm, align=center, fill=blue!8},
 rewrite-technology/.style={draw, rounded corners=2pt, minimum width=3.5cm,
  minimum height=1.0cm, align=center, fill=orange!9},
 rewrite-flow/.style={->, semithick},
 rewrite-dynamics/.style={->, semithick, dashed}
}

\subsection{Population and labor-augmenting technology}
\label{subsec:rewrite-population}

Population grows at the exogenous rate $n\geq0$:
\begin{equation}
 \dot N=nN,
 \qquad n\geq0.
 \label{eq:rewrite-population}
\end{equation}
Each person supplies one unit of homogeneous labor inelastically. The aggregate labor endowment is therefore $N$. Because autonomous research uses no human labor, all labor is employed in final production and equilibrium will impose $L=N$.

The variable $A$ measures the level of labor-augmenting technology. One unit of labor supplies $A$ efficiency units, so effective labor is $AL$. Labor-augmenting technology evolves according to
\begin{equation}
 \dot A=\gamma A,
 \qquad \gamma>0.
 \label{eq:rewrite-labor-productivity}
\end{equation}
%For any positive variable $V$, I write $g_V\equiv\dot V/V$. Thus $g_A=\gamma$.

\subsection{Household}
\label{subsec:rewrite-household}

The representative household owns physical capital and the AI developer. Let \(\Pi\) denote the AI developer's net profit, which accrues to the household. The household takes the paths of the net interest rate \(r\), the wage \(w\) in final-good units and \(\Pi\) as given. For a predetermined capital stock $K_0>0$, it chooses aggregate consumption $C$ to solve
\begin{equation}
 \max_{\{C_t\}_{t\geq0}}
 \int_0^\infty e^{-\rho t}N_t\log(C_t/N_t)\dd t,
 \qquad \rho>n,
 \label{eq:rewrite-household-problem}
\end{equation}
subject to
\begin{equation}
 \dot K=rK+wN+\Pi-C,\qquad K\geq0.
 \label{eq:rewrite-household-budget}
\end{equation}
The variable $r$ is the net real interest rate: depreciation has already been deducted from the rental return to capital. The household's Euler condition is
\begin{equation}
 \frac{\dot C}{C}=n+r-\rho.
 \label{eq:rewrite-euler}
\end{equation}
Its transversality condition is
\begin{equation}
 \lim_{t\to\infty}e^{-\rho t}\frac{N_tK_t}{C_t}=0.
 \label{eq:rewrite-household-tvc}
\end{equation}
%These conditions use only objects in the household's problem. Production-side conditions determining $r,w,$ and \(\Pi\) are introduced below.

\subsection{Final-good production}
\label{subsec:rewrite-final-production}

Competitive firms produce the final good from physical capital $K$, effective labor $AL$, and effective AI production services $X$. Let $0<\alpha<1$ denote the capital share, $\delta\geq0$ the depreciation rate, and $\sigma>0$ the elasticity of substitution between effective labor and effective AI production services. The parameters $\omega_L$ and $\omega_X$ denote the CES weights on effective labor and effective AI production services, respectively. Production is
\begin{align}
 Y&=K^\alpha Z^{1-\alpha},
 \label{eq:rewrite-output}\\
 Z&=\left[\omega_L(AL)^{(\sigma-1)/\sigma}
 +\omega_XX^{(\sigma-1)/\sigma}\right]^{
 \sigma/(\sigma-1)}.
 \label{eq:rewrite-ces}
\end{align}
The composite input $Z$ combines effective labor and AI production services.
%At the continuous unit-elasticity limit,
%$Z=(AL)^{\omega_L}X^{\omega_X}$.

%To illustrate this production technology, consider a researcher writing an academic paper. The researcher contributes time and expertise to formulate questions, evaluate arguments, and interpret results. AI assists with derivations, programming, and drafting, while physical capital, such as a computer, runs numerical calculations and produces figures. AI can replace some of the work the researcher would otherwise do, but it can also make the researcher's expertise more productive by helping test conjectures or communicate arguments more clearly. The researcher may revise an argument drafted by AI or ask AI to check an argument written by the researcher. A numerical result may similarly combine a question formulated by the researcher, code developed with AI assistance, and calculations executed on the computer.

Jobs comprise multiple tasks, and AI may substitute for human input in some tasks while complementing it in others \citep{autor2013,acemoglurestrepo2019}.
%An early model of mechanization is \citet{zeira1998}, in which producers can use more capital to reduce the labor required to produce individual intermediate goods. 
%I model labor and AI services as aggregate inputs, not as alternative ways of performing a single task. 
\citet{jonestonetti2026} show how a production process with many tasks can admit a CES aggregate representation. \citet{erdiletal2025gate} distinguish automating additional tasks (the extensive margin) from increasing the services supplied in tasks that AI
can already perform (the intensive margin). I focus on expansion of aggregate AI services without explicitly tracking the set of automatable tasks.

Let \(p_X\) denote the price of one unit of effective AI production services. Taking the input prices \(r+\delta,w,\) and \(p_X\) as given, a final-good firm solves
\begin{equation}
 \max_{K,L,X\geq0}
 \left\{K^\alpha Z(AL,X)^{1-\alpha}-(r+\delta)K-wL-p_XX\right\}.
 \label{eq:rewrite-final-firm-problem}
\end{equation}
Define $s_X$ as the elasticity of $Z$ with respect to $X$:
\begin{equation}
 s_X=\frac{\omega_XX^{(\sigma-1)/\sigma}}
 {\omega_L(AL)^{(\sigma-1)/\sigma}
 +\omega_XX^{(\sigma-1)/\sigma}}.
 \label{eq:rewrite-sx}
\end{equation}
The first-order conditions are
\begin{align}
 r&=\alpha\frac{Y}{K}-\delta,
 \label{eq:rewrite-capital-return}\\
 w&=(1-\alpha)(1-s_X)\frac{Y}{L},
 \label{eq:rewrite-wage}\\
 p_X&=(1-\alpha)s_X\frac{Y}{X}.
 \label{eq:rewrite-ai-price}
\end{align}
Thus $(1-\alpha)s_X$ is AI revenue as a share of output, while
$(1-\alpha)(1-s_X)$ is labor's income share.

Constant returns and these conditions imply
\begin{equation}
 Y=(r+\delta)K+wL+p_XX.
 \label{eq:rewrite-final-zero-profit}
\end{equation}

Solving the final-good firm's AI condition for its price yields the inverse
demand schedule $p_X(X)$. %This notation conditions on the current values of $K,A,$ and $L$; I display those arguments only when the time dependence matters.
The absolute elasticity of inverse demand, holding $K$, $A$, and $L$ fixed, is
\begin{equation}
 e_X\equiv-
 \frac{\partial\ln p_X(X)}{\partial\ln X}
 =\frac{1-s_X}{\sigma}+\alpha s_X.
 \label{eq:rewrite-inverse-demand-elasticity}
\end{equation}
This demand schedule is an input to the developer's problem.

\subsection{The AI developer}
\label{subsec:rewrite-developer}

\subsubsection{Effective AI production and research services}
\label{subsubsec:rewrite-ai-services}

I use \emph{compute} to denote the computational resources used by
AI systems \citep{sastryetal2024}. %The variables $U$ and $M$ are rates of compute use rather than stocks of processors or data centers.
Compute can be measured in
floating-point operations, with flows measured in operations per unit
of time. I assume that each unit of compute requires one unit of the
final good, which could otherwise be consumed or invested in physical
capital.

Inference compute $U$ supplies AI services for current final-good
production, while research compute $M$ is used to develop better AI.
\citet{korinekmckelvey2026} make a related distinction between compute
used for inference and compute used for R\&D and training. Research
compute includes both training and running AI systems that conduct
research \citep{davidsonetal2026}. The distinction between $U$ and $M$
therefore concerns their economic use, not a strict separation between
inference and training.

%In practice, AI research also uses human researchers, data, and complementary infrastructure. The benchmark does not treat $M$ as a complete measure of R\&D. It isolates its computational component and assumes that human research is no longer the binding input.

AI efficiency $B$ measures the services obtained per unit of compute.
I assume that the same $B$ augments both activities, producing flows
$BU$ and $BM$ of AI production and research services, respectively.
A rise in $B$ represents algorithmic improvement rather than an
expansion of physical computing capacity. This interpretation parallels
measures of algorithmic efficiency based on the compute required to
achieve a given level of performance
\citep{hernandezbrown2020,erdilbesiroglu2022}.
AI production services therefore satisfy
\begin{equation}
 X=BU.
 \label{eq:rewrite-ai-service}
\end{equation}

Figure~\ref{fig:rewrite-model-flow} summarizes the production block.

\begin{figure}[H]
 \centering
 \resizebox{0.84\textwidth}{!}{%
 \begin{tikzpicture}[>=Latex,font=\small]
  \node[rewrite-state] (A) at (-7.5,2.0) {\(A\)\\labor-augmenting\\technology};
  \node[rewrite-input] (L) at (-7.5,0.7) {\(L\)\\production labor};
  \node[rewrite-technology] (AL) at (-4.0,1.35) {\(AL\)\\effective labor};
  \node[rewrite-state] (B) at (-7.5,-0.8) {\(B\)\\AI efficiency};
  \node[rewrite-input] (U) at (-7.5,-2.1) {\(U\)\\inference compute};
  \node[rewrite-technology] (X) at (-4.0,-1.45) {\(X=BU\)\\AI production services};
  \node[rewrite-technology] (Z) at (-0.2,0.0)
   {\(Z=\operatorname{CES}(AL,X)\)\\elasticity \(\sigma\)};
  \node[rewrite-state] (K) at (-0.2,2.2) {\(K\)\\physical capital};
  \node[rewrite-technology] (Y) at (3.5,1.2)
   {\(Y=K^\alpha Z^{1-\alpha}\)\\final output};

  \draw[rewrite-flow] (A) -- (AL);
  \draw[rewrite-flow] (L) -- (AL);
  \draw[rewrite-flow] (B) -- (X);
  \draw[rewrite-flow] (U) -- (X);
  \draw[rewrite-flow] (AL) -- (Z);
  \draw[rewrite-flow] (X) -- (Z);
  \draw[rewrite-flow] (Z) -- (Y);
  \draw[rewrite-flow] (K) -- (Y);
 \end{tikzpicture}%
 }
 \caption{Production of final output and AI production services. AI efficiency $B$ augments inference compute $U$ to produce $X=BU$, the flow of AI production services. These services combine with effective labor $AL$ to produce $Z$, which combines with capital $K$ to produce $Y$.}
 \label{fig:rewrite-model-flow}
\end{figure}

\subsubsection{RSI technology}
\label{subsubsec:rewrite-capability}

Given initial AI efficiency \(B_0>0\) and an upper bound
\(\overline B\in(B_0,\infty]\), AI efficiency evolves according to the RSI technology
\begin{equation}
 \dot B=\chi(BM)^\eta\psi(B),
 \qquad
 \psi(B)=
  1-\frac{B}{\overline B},
 \qquad \chi>0,\quad0<\eta<1.
 \label{eq:rewrite-capability-law}
\end{equation}
Figure~\ref{fig:rewrite-capability-flow} summarizes the RSI mechanism.
\begin{figure}[H]
 \centering
 \resizebox{0.72\textwidth}{!}{%
 \begin{tikzpicture}[>=Latex,font=\small]
  \node[rewrite-input] (M) at (-5.9,0.9) {\(M\)\\research compute};
  \node[rewrite-state] (B) at (-5.9,-0.9) {\(B\)\\AI efficiency};
  \node[rewrite-technology] (BM) at (-1.9,0.9)
   {\(BM\)\\AI research services};
  \node[rewrite-technology] (psi) at (-1.9,-0.9)
   {\(\psi(B)\)\\remaining room};
  \node[rewrite-technology] (Bd) at (2.2,0.0)
   {\(\dot B=\chi(BM)^\eta\psi(B)\)\\efficiency improvement};

  \draw[rewrite-flow] (B) -- (BM);
  \draw[rewrite-flow] (M) -- (BM);
  \draw[rewrite-flow] (B) -- (psi);
  \draw[rewrite-flow] (BM) -- (Bd);
  \draw[rewrite-flow] (psi) -- (Bd);
  \draw[rewrite-dynamics] (Bd.east) -- (4.4,0.0) -- (4.4,-2.0) --
    (-5.9,-2.0) -- (B.south);
 \end{tikzpicture}%
 }
 \caption{Recursive self-improvement in AI efficiency. AI efficiency $B$ augments
 research compute $M$ to produce $BM$, the flow of AI research services. These services and the remaining
 room for improvement $\psi(B)$ determine \(\dot B\), which increases $B$.}
 \label{fig:rewrite-capability-flow}
\end{figure}
The parameter \(\chi\) shifts research productivity, while \(\eta\) is the elasticity of \(\dot B\) with respect to research compute, holding $B$ fixed. Equation~\eqref{eq:rewrite-capability-law} captures RSI: a higher $B$ produces more AI research services $BM$ from a given amount of research compute $M$, and these services help improve $B$ further. This is the AI-research feedback studied by \citet{davidsonetal2026}. The term \(\psi(B)\) reduces the remaining room for improvement as AI efficiency approaches \(\overline B\). A related approach appears in \citet{erdiletal2025gate}, who also reduce research productivity as hardware and software efficiency approach their respective upper bounds. The uncapped specification is recovered at \(\overline B=\infty\).

Evidence that ideas become harder to find motivates allowing research
productivity to decline as technology advances. \citet{bloometal2020} document that research effort has risen substantially while research productivity has fallen across several industries, products, and firms. %Following thisinterpretation, \citet{davidsonetal2026} use their estimates to discipline diminishing returns in idea production. 
This evidence motivates diminishing returns, but it does not establish a literal finite technological ceiling. The term $1-B/\overline B$ is a deliberately stronger reduced-form specification. Without the upper bound, recursive improvement may lead to the absence of an equilibrium, as shown in Section~\ref{subsec:rewrite-uncapped}. %instead produce a finite-time singularity: one or more endogenous quantities---such as AI efficiency, output, capital, or research compute---become unbounded at a finite date, so the trajectory cannot be continued as a path of finite allocations. %Section~\ref{subsec:rewrite-uncapped} explains why the uncapped economy requires a separate equilibrium analysis. %AI efficiency may approach the upper bound asymptotically but cannot cross it in finite time. %Appendix~\ref{app:rewrite-cap-invariance} gives the derivation.

\subsubsection{Supply of effective AI production services}
\label{subsubsec:rewrite-market-structure}

An integrated AI developer has exclusive ownership of $B$. It conducts research
and sells effective AI production services. The developer sells $X$ and purchases the $U=X/B$ units of inference compute needed to supply those services. Because $X$ does not enter the RSI technology in equation~\eqref{eq:rewrite-capability-law}
directly, its within-date choice can be solved conditional on $B$.
Define optimized operating profit before research expenditure as
\begin{equation}
 \pi(B)=\max_{X\geq0}
 \left\{p_X(X)X-\frac{X}{B}\right\}.
 \label{eq:rewrite-service-choice}
\end{equation}
The notation \(\pi(B)\) leaves its dependence on current \(K,A,\) and \(L\) implicit.
The first-order condition equates marginal revenue to the compute cost of one additional service unit:
\begin{equation}
 p_X(1-e_X)=\frac{1}{B}.
 \label{eq:rewrite-monopoly-foc}
\end{equation}

Lemma~\ref{lem:rewrite-monopoly-choice} in
Appendix~\ref{app:rewrite-proofs} shows that
problem~\eqref{eq:rewrite-service-choice} has a unique, strictly positive
global maximizer characterized by condition~\eqref{eq:rewrite-monopoly-foc}.
Thus, the developer charges the unit cost $1/B$ multiplied by the monopoly
markup $1/(1-e_X)$.
%At \(\sigma=1\), define
%\(\beta=(1-\alpha)\omega_X\). The solution is explicit:
%\begin{equation}
% X^*=\left[\beta^2BK^\alpha
% (AL)^{(1-\alpha)\omega_L}\right]^{1/(1-\beta)}.
% \label{eq:rewrite-monopoly-cd-solution}
%\end{equation}
The envelope theorem gives
\begin{equation}
 \pi_B=\frac{X}{B^2}=\frac{U}{B}.
 \label{eq:rewrite-profit-envelope}
\end{equation}
%The $B^2$ term does not represent quadratic returns to AI efficiency. It results from differentiating the inference cost $X/B$:
At a fixed $X$, a marginal increase in AI efficiency saves $U/B$ units of the consumption good.

AI efficiency is nonrival in use. Once $B$ has been developed, using it to
augment one unit of inference or research compute does not reduce its ability to augment other units. This does not make effective AI services costless: $BU$ and $BM$ require the rival and costly compute inputs $U$ and $M$. Nonrivalry also does not rule out excludability. An owner can restrict access to model weights, software, or AI services.

To see the resulting appropriability problem, suppose that competitive
service providers had free access to $B$. Free entry would imply no operating rent. A developer that paid for $M$ would improve a common AI-efficiency index $B$ that competing providers could use without compensating it. The developer would bear the research cost but could not appropriate the resulting benefits. %The problem is therefor one of research incentives, not the physical availability of compute.
%I therefore assign exclusive ownership of $B$ to an integrated developer that conducts research and sells effective AI production services.
Exclusive ownership allows the developer to capture part of the value created by a higher $B$ and makes the return to research appropriable, as in the standard endogenous-growth treatment of nonrival ideas \citep{romer1990}.
That being said, monopoly is an institutional assumption, not a technological implication of nonrivalry.\footnote{Licensing, oligopoly, public funding, or open-source provision with research subsidies could also support investment in $B$ \citep{korinekvipra2025,gans2024,demireretal2025}. I also abstract from the conflict between general-purpose AI suppliers and application developers studied by \citet{liuwan2026}.}

\subsubsection{The dynamic problem}
\label{subsubsec:rewrite-dynamic-developer}

%Taking the paths of $r,K,A,$ and $L$ as given, it solves
%\begin{equation}
% \max_{\{X_t,M_t\}_{t\geq0}}
% \int_0^\infty
% \exp\left\{-\int_0^t r_s\dd s\right\}
% \left[p_X(X_t;K_t,A_t,L_t)X_t-\frac{X_t}{B_t}-M_t\right]\dd t,
% \label{eq:rewrite-developer-problem}
%\end{equation}
%subject to \eqref{eq:rewrite-capability-law} and $B(0)=B_0$. Its distributed
%profit is
%\begin{equation}
% \Pi=p_XX-\frac{X}{B}-M=p_XX-U-M.
% \label{eq:rewrite-developer-profit}
%\end{equation}
%Profit can be negative at a date because the model imposes no separatecash-flow constraint. A negative \(\Pi\) represents an equity injection from the household that owns the developer.
%\subsubsection{Pointwise service choice}
%\label{subsubsec:rewrite-static-monopoly}
%Sales $X$ do not enter the capability law or any other intertemporal
%constraint. Conditional on $B,K,A,$ and $L$, the developer can therefore
%choose $X$ within each date. Define optimized operating profit before
%research expenditure as
%\begin{equation}
% \pi(B;K,A,L)=\max_{X\geq0}
% \left\{p_X(X)X-\frac{X}{B}\right\}.
% \label{eq:rewrite-service-choice}
%\end{equation}
%The first-order condition equates marginal revenue to the compute cost of one additional service unit:
%\begin{equation}
% p_X(1-e_X)=\frac{1}{B}.
% \label{eq:rewrite-monopoly-foc}
%\end{equation}
%\begin{proposition}[Pointwise monopoly choice]
%\label{prop:rewrite-monopoly-choice}
%Suppose $K,A,L,B>0$, $0<\alpha<1$, $\sigma>0$, and
%\(\omega_L,\omega_X\in(0,1)\). Problem \eqref{eq:rewrite-service-choice} has a
%unique global maximizer $X^*>0$. It is the unique solution to
%\eqref{eq:rewrite-monopoly-foc} and satisfies
%\begin{equation}
% e_X(1-e_X)
% +(\alpha-\sigma^{-1})(1-\sigma^{-1})s_X(1-s_X)>0.
% \label{eq:rewrite-monopoly-soc}
%\end{equation}
%\end{proposition}
%The monopoly price is therefore
%\begin{equation}
% p_X=\frac{1/B}{1-e_X}.
% \label{eq:rewrite-markup}
%\end{equation}
%At \(\sigma=1\), define
%\(\beta=(1-\alpha)\omega_X\). The solution is explicit:
%\begin{equation}
% X^*=\left[\beta^2BK^\alpha
% (AL)^{(1-\alpha)\omega_L}\right]^{1/(1-\beta)}.
% \label{eq:rewrite-monopoly-cd-solution}
%\end{equation}
%The envelope theorem gives
%\begin{equation}
% \pi_B=\frac{X}{B^2}=\frac{U}{B}.
% \label{eq:rewrite-profit-envelope}
%\end{equation}
%The $B^2$ term does not represent quadratic returns to AI efficiency. It results from differentiating the inference cost $X/B$: at a fixed $X$, a marginal increase in AI efficiency saves $X/B^2$ units of compute.
After the within-date service choice, the developer's net profit is
$\Pi_t=\pi(B_t)-M_t$: operating profit less research expenditure.
Taking the paths of $r,K,A,$ and $L$ as given and substituting the static
solution to problem~\eqref{eq:rewrite-service-choice}, the developer chooses
the path of research compute $M$ to increase AI efficiency $B$:
\begin{equation}
\max_{\{M_t\}_{t\geq0}}
 \int_0^\infty
 \exp\left\{-\int_0^t r_s\dd s\right\}
 [\pi(B_t)-M_t]\dd t,
 \label{eq:rewrite-developer-reduced}
\end{equation}
subject to the RSI technology given by \eqref{eq:rewrite-capability-law} and taking $B_0$ as given. %This reduction does not replace the dynamic problem with a static problem. It solves the within-date choice $X$ inside the original intertemporal problem and leaves $M$ as the dynamic control.
%Admissible research paths have $M_t\geq0$ and finite total expenditure on every finite interval. The objective is the present value of net profit, not profit at a single date.

%Profit can be negative at a date because the model imposes no separate cash-flow constraint. A negative \(\Pi\) represents an equity injection from the household that owns the developer.

Let $q$ be the current-value shadow price of AI efficiency, measured in units of
the current final good per unit of $B$. An interior choice of research compute $M$ satisfies
\begin{equation}
 q\chi\eta B(BM)^{\eta-1}\psi(B)=1,
 \label{eq:rewrite-research-foc}
\end{equation}
%or, equivalently,
%\begin{equation}
% M=\left[q\chi\eta B^\eta\psi(B)
% \right]^{1/(1-\eta)}.
% \label{eq:rewrite-research-demand}
%\end{equation}
Equation~\eqref{eq:rewrite-research-foc} equates the marginal benefit of research compute to its unit cost. Holding current $B$ fixed, an additional unit of $M$ raises the rate of improvement in AI efficiency by \(\chi\eta B(BM)^{\eta-1}\psi(B)\).
%The factor $B$ appears because raw compute $M$ produces $BM$ effective AI research services.
Multiplication by $q$ values the resulting efficiency gain, accounting for future profit gains and further improvements in AI efficiency. %The right-hand side is one because one unit of compute costs one unit of the final good. At an interior optimum, an additional unit of research therefore creates exactly enough value to cover its cost.

The current-value costate equation is
\begin{equation}
 \frac{\dot q}{q}=r-\frac{X}{qB^2}-\eta\frac{\dot B}{B}
 -\dot B\frac{\psi'(B)}{\psi(B)}.
 \label{eq:rewrite-costate}
\end{equation}
%For a finite upper bound,\(\psi'/\psi=-1/(\overline B-B)\); for the uncapped economy, the derivative is zero.
Rearranging gives the return equation
\begin{equation}
 rq=\frac{U}{B}+q\eta\frac{\dot B}{B}
 +q\dot B\frac{\psi'(B)}{\psi(B)}+\dot q.
 \label{eq:rewrite-capability-return}
\end{equation}
The first term on the right-hand side of Equation~\eqref{eq:rewrite-capability-return} is the inference-compute saving from greater AI efficiency. The
second is the research-productivity gain from RSI. With an upper bound, the third
term is negative because a higher $B$ leaves less room for further improvement. The final term is the capital gain on AI efficiency. Together these
returns must equal the market return $rq$.

The developer's transversality condition is
\begin{equation}
 \lim_{t\to\infty}
 \exp\left\{-\int_0^t r_s\dd s\right\}q_tB_t=0.
 \label{eq:rewrite-developer-tvc}
\end{equation}
It rules out a path that leaves positive discounted value in the AI-efficiency state.%It does not require $q$, $B$, or their undiscounted product to converge to zero.

The preceding conditions are necessary for an interior optimum but do not
establish global optimality.
Lemma~\ref{lem:rewrite-developer-verification} in
Appendix~\ref{app:rewrite-proofs} gives sufficient conditions for a trajectory
of research expenditure to maximize the present value of the developer's net profit in
problem~\eqref{eq:rewrite-developer-reduced}.

%\begin{proposition}[Research is interior]
%\label{prop:rewrite-research-interior}
%Suppose $0<\eta<1$ and the developer has a differentiable, finite-valued optimum with $B>0$ and positive demand for effective AI production services. Then $q>0$ and $M>0$ at every finite date. Research may converge to zero asymptotically, but the developer cannot stop improvements in AI efficiency by choosing $M=0$ at a finite date.
%\end{proposition}

\subsection{Equilibrium}
\label{subsec:rewrite-equilibrium}

Capital-market clearing
identifies household capital with the capital rented by final-good firms.
Labor-market clearing gives $L=N$. AI-service market clearing identifies the
developer's sales with the final sector's purchases. Combining the household
budget, final-good zero profit, developer profit, and market clearing yields
the aggregate resource constraint
\begin{equation}
 C+\dot K+\delta K+U+M=Y.
 \label{eq:rewrite-resource}
\end{equation}

\begin{definition}[Equilibrium]
\label{def:rewrite-equilibrium}
For exogenous initial values $A_0,N_0,K_0,B_0>0$, an equilibrium is an
infinite-horizon trajectory
\begin{equation}
 \left\{K_t,B_t,C_t,q_t,L_t,U_t,M_t,X_t,Z_t,Y_t,r_t,w_t,p_{X,t},
 s_{X,t},e_{X,t},\Pi_t\right\}_{t\geq0}
 \label{eq:rewrite-equilibrium-trajectory}
\end{equation}
such that the developer solves problem~\eqref{eq:rewrite-developer-reduced}
and the following conditions hold:

%The exogenous paths and predetermined initial conditions are
%\begin{subequations}
%\label{eq:rewrite-equilibrium-exogenous}
%\begin{align}
% A_t&=A_0e^{\gamma t},\\
% N_t&=N_0e^{nt},%\\
% %K(0)&=K_0, & B(0)&=B_0.
%\end{align}
%\end{subequations}
\begin{enumerate}
\item The intratemporal allocation and prices satisfy
\begin{subequations}
\label{eq:rewrite-equilibrium-static}
\begin{align}
 L&=N,
 \label{eq:rewrite-eq-labor}\\
 Z&=\left[\omega_L(AL)^{(\sigma-1)/\sigma}
 +\omega_XX^{(\sigma-1)/\sigma}\right]^{
 \sigma/(\sigma-1)},
 \label{eq:rewrite-eq-z}\\
 Y&=K^\alpha Z^{1-\alpha},
 \label{eq:rewrite-eq-y}\\
 s_X&=\frac{\omega_XX^{(\sigma-1)/\sigma}}
 {\omega_L(AL)^{(\sigma-1)/\sigma}
 +\omega_XX^{(\sigma-1)/\sigma}},
 \label{eq:rewrite-eq-sx}\\
 e_X&=\frac{1-s_X}{\sigma}+\alpha s_X,
 \label{eq:rewrite-eq-ex}\\
 r&=\alpha\frac{Y}{K}-\delta,
 \label{eq:rewrite-eq-r}\\
 w&=(1-\alpha)(1-s_X)\frac{Y}{L},
 \label{eq:rewrite-eq-w}\\
 p_X&=(1-\alpha)s_X\frac{Y}{X},
 \label{eq:rewrite-eq-px}\\
 \frac{1}{B}&=p_X(1-e_X),
 \label{eq:rewrite-eq-monopoly}\\
 X&=BU,
 \label{eq:rewrite-eq-service}\\
 1&=q\chi\eta B(BM)^{\eta-1}\psi(B),
 \label{eq:rewrite-eq-research}\\
 \Pi&=p_XX-U-M.
 \label{eq:rewrite-eq-profit}
\end{align}
\end{subequations}
\item The four dynamic equations are
\begin{subequations}
\label{eq:rewrite-equilibrium-dynamics}
\begin{align}
 \dot K&=Y-C-U-M-\delta K,
 \label{eq:rewrite-eq-kdot}\\
 \dot B&=\chi(BM)^\eta\psi(B),
 \label{eq:rewrite-eq-bdot}\\
 \frac{\dot C}{C}&=n+r-\rho,
 \label{eq:rewrite-eq-cdot}\\
 \frac{\dot q}{q}&=r-\frac{X}{qB^2}-\eta\frac{\dot B}{B}
 -\dot B\frac{\psi'(B)}{\psi(B)}.
 \label{eq:rewrite-eq-qdot}
\end{align}
\end{subequations}
\item The two transversality conditions are
\begin{subequations}
\label{eq:rewrite-equilibrium-tvcs}
\begin{align}
 \lim_{t\to\infty}e^{-\rho t}\frac{N_tK_t}{C_t}&=0,
 \label{eq:rewrite-eq-household-tvc}\\
 \lim_{t\to\infty}
 \exp\left\{-\int_0^t r_s\dd s\right\}q_tB_t&=0.
 \label{eq:rewrite-eq-developer-tvc}
\end{align}
\end{subequations}
\end{enumerate}
%At \(\sigma=1\), equations \eqref{eq:rewrite-eq-z} and \eqref{eq:rewrite-eq-sx} are understood at their continuous limits:
%$Z=(AL)^{\omega_L}X^{\omega_X}$ and $s_X=\omega_X$. If \(\overline B<\infty\), the path satisfies $B_t<\overline B$ at every finite date. Finally, t
%The household and developer solve their stated optimization problems, final-good firms maximize profit, and all markets clear. The initial values \(C_t\) and \(q_t\) are endogenous jump variables selected to satisfy the two transversality conditions.
\end{definition}

%The definition does not impose balanced growth. A balanced-growth path is one possible equilibrium trajectory, not the meaning of equilibrium. Nor are the dated equations alone enough: a numerical path is admitted as an equilibrium only if it also satisfies both transversality conditions and the global optimality requirements in Definition \ref{def:rewrite-equilibrium}. 
\begin{remark}[The economy without AI]
\label{rem:rewrite-no-ai}
With $\omega_X=0$ and $\omega_L=1$, the model reduces to the
Ramsey--Cass--Koopmans economy with logarithmic utility and exogenous
labor-augmenting technological progress.
\end{remark}

The system contains no separate household budget or final-good zero-profit
equation because these are redundant after market clearing. The consolidated
income decomposition follows from the equations in the definition:
\begin{equation}
 1=\frac{(r+\delta)K}{Y}+\frac{wL}{Y}
 +\frac{\Pi}{Y}+\frac{U}{Y}+\frac{M}{Y}.
 \label{eq:rewrite-income-decomposition}
\end{equation}
It records how final output is divided among gross capital income, labor income,
developer profit, inference compute, and research compute. %It is an accounting consequence of equilibrium, not an additional equilibrium condition.
